% ============================================================================
% sec-status.tex — Formalization status, limitations, and open problems.
% Label prefix: sta: (all labels owned by this file).
% ============================================================================
\section{Formalization status, limitations, and open problems}\label{sec:open}

\subsection{The Lean formalization, item by item}\label{sta:subsec-leanitems}

Table~\ref{sta:tab-leanitems} lists every Lean declaration this paper
cites that is free of \texttt{sorryAx}, with its exact statement, source
anchor and axiom profile; every entry of the axiom column was audited by
\texttt{\#print axioms} against the built library, and the audit was
re-run against the current sources (identical output; the audited files
byte-identical to the review snapshot; Section~\ref{com:subsec-lean}).
Declarations whose proofs consume \texttt{native\_decide} carry the
kernel-reflection axiom \texttt{Lean.ofReduceBool} explicitly and are
quoted as ``machine-checked by kernel reduction''
(Section~\ref{com:subsec-lean}); they are not equated with pure proofs.

\begin{table}[ht]
\centering\small
\caption{The nine \texttt{sorryAx}-free Lean items cited by this paper.
Axiom profiles: $[\mathrm{S}]:=$ \texttt{[propext, Classical.choice,
Quot.sound]} (standard), $[\mathrm{P},\mathrm{Q}]:=$
\texttt{[propext, Quot.sound]}, and $\mathrm{nd}:=$ the
\texttt{Lean.ofReduceBool} kernel-reflection axiom instance. Line anchors
refer to the released sources where the file is released;
\texttt{CompleteOdd.lean} and \texttt{Cert8.lean} are not part of the
released snapshot and their anchors refer to the audited development
snapshot of the project.}
\label{sta:tab-leanitems}
\begin{tabular}{p{3.55cm} p{4.55cm} p{2.7cm} p{2.0cm}}
\hline
Declaration & Statement (exactly as proved) & Anchor & Axioms\\
\hline
\texttt{pNumK4n\_odd\_all} & for every odd $n\ge3$:
$\pn(K_{4,n})=(n+1)/2$; hypotheses $3\le n$, $n$ odd & \texttt{T273.lean}
:976--978 & $[\mathrm{S}]$\\
\texttt{cor\_pNumK4n\_5/13/19} & $\pn(K_{4,5})=3$,
$\pn(K_{4,13})=7$, $\pn(K_{4,19})=10$ & \texttt{T273.lean} :981--987 &
$[\mathrm{S}]$\\
\texttt{pCompleteOdd} & for $k\ge1$: $\pn(K_{2k+1})=k+1$; hypothesis
$k\ne0$ & \texttt{CompleteOdd.lean} :1590--1596 & $[\mathrm{S}]$\\
\texttt{thmA2663\_upgraded} & under hypotheses
$\mathrm{hd2}$ (target side has the prescribed size) and
$\mathrm{val}\le1$: $g_{13}\le d_2+1$ --- the $\mathrm{val}\le1$
hypothesis is \emph{not} proved by the theorem &
\texttt{ThmA2663.lean} :197--200 & $[\mathrm{P},\mathrm{Q}]$\\
\texttt{case2\_struct\_alpha\_iff} & under six explicit hypotheses:
$C_1((t_1+g)\bmod N)=\texttt{true}\iff t_1=s$ & \texttt{Struct2663.lean}
:260--265 & $[\mathrm{P},\mathrm{Q}]$\\
\texttt{ledger6\_ok} & $\mathrm{ledger}_6=(5{,}724;\,2{,}613;\,190;\,398;\,588)$
(N=6 five-column global account) & \texttt{Cert2663.lean} :113 &
$[\mathrm{P},\mathrm{nd}]$\\
\texttt{ledger8\_ok} & $\mathrm{ledger}_8=(701{,}360;\,285{,}372;\,20{,}308;\,10{,}452;\,30{,}760)$
(N=8; dual-route, second Python implementation pinned) &
\texttt{Cert8.lean} :52 & $[\mathrm{P},\mathrm{nd}]$\\
\texttt{l2b\_toplevel\_certificate} & the $79{,}200$-point packing-table
certificate ($120$ orbits $\times$ $330$ demands $\times$ two sides),
pointwise iff & \texttt{Pi3L2b.lean} :248--249 & $[\mathrm{P},\mathrm{nd}]$\\
\texttt{Pi3Uni} (nine theorems) & downset tables (both sides); uniform
layer; the no-threshold witness refuting any ``feasible $\iff
f(\Sigma d)$'' law; top-only failures ($1{,}761$); top-failure set an
upset; orbit counts $110/2$ & \texttt{Pi3Uni.lean} :57--115 &
$[\mathrm{P},\mathrm{nd}]$\\
\hline
\end{tabular}
\end{table}

\emph{Not formalized.} Three boundary entries complete the picture.
Theorem~\ref{thm:conchain} (odd clique chains) and its engine
Theorem~\ref{thm:conL0} have \emph{no} Lean statement: they are paper
proofs, and the machine anchors cover the underlying decompositions, not
the chaining argument (Section~\ref{subsec:conformal}). The Fan-type chain
(Theorems~\ref{thm:claAstar},~\ref{thm:lasep}) carries
\texttt{sorryAx}-dependent declarations --- formalization in progress, and
never quoted as machine-verified (Remark~\ref{rem:claformal}). The
research lemma \texttt{fan6\_nf} is never cited: its statement was
refuted by a machine counterexample and it survives in the sources as a
record only.

\emph{The four-domain account layer and its honest debt list.} The
termination bound's four-domain witness accounts
($588/30{,}760/634{,}960/7{,}798{,}903$ instances at $N=6,8,10,12$;
Section~\ref{com:subsec-lean}) are kernel-checked equalities at the level
of the domain counts and of the two certified ledgers; the theorem body
itself is a paper proof with no Lean statement. Five residuals are
registered against the account layer, none claimed closed:
(R1) the fast$\equiv$slow enumeration equivalence has no general-$N$
theorem (it is pinned by double-scale agreement at the $N=6/8$ anchors and
across all bucket tables); (R2) the full witness tables at $N=10$
($12{,}211{,}196$ rows) and $N=12$ ($313{,}315{,}956$ rows) are not
certified --- bucket-first sampling and the aggregate bucketed ledgers
cover them; (R3) the trailing four digits of the $N=10$/$N=12$ tuples have
no kernel-level \emph{third} implementation (the Lean pins and the Python
second implementation make two); (R4) the $h\bmod 2g$ keys bridge is
proved at zero axioms for its instances, but its general-$N$ derivation is
open; (R5) the declaration-listing verifier is a text-level instrument
whose statement-level boundary is carried as such.

\subsection{Limitations}\label{sta:subsec-limitations}

\subsubsection{The budget-conjecture residuals}\label{sta:subsubsec-kres}

\paragraph{The degree-two cut-vertex branch.} The degree-two layer of the
budget conjecture assembles through a case analysis whose cut-vertex
branch ($G-x$ disconnected) is closed at the \emph{computational} level
only: $1{,}870$ cut-vertex even-degree instances over the base domain of
Table~\ref{com:tab-kdomains} ($492$ with $d=2$, $1{,}165$ with $d=4$,
$213$ with $d=6$) show zero violations, and a dedicated connected
$n\le7$ check covers all $106$ degree-$2$ cut-vertex instances, again
with zero violations. The theorization is pending, and the drafted
splicing argument has a genuine gap: its parity arithmetic assumes
$\lceil n_1/2\rceil+\lceil n_2/2\rceil=n/2$ for the two components of
$G-x$, which holds only when $n$ is even (components of opposite
parity). For odd $n$ the components share parity: the same-odd split
gives $(n+1)/2$ and the same-even split gives $(n-1)/2$, so the drafted
equality fails by one --- the missing parity case must be supplied
before the branch is a theorem. The unconditional ceiling survives
unchanged: $\lceil n_1/2\rceil+\lceil n_2/2\rceil\le\lceil n/2\rceil$
for all parity splits, verified exhaustively for $n\le30$.

\paragraph{The degree-$\ge4$ layer.} The general degree-$\ge4$ layer of
the budget conjecture is \emph{not} discharged. The exact criterion
(Section~\ref{int:subsec-budget}) characterizes one construction tool ---
the degree-insertion operations --- and the alternative-construction
theory for the $208$ tool-dead instances of
Table~\ref{com:tab-kdomains} ($180$ of $208$ cut-vertex; only $7$ with
the reduced case $\sigma=0$) is not written; the enumeration is confined
to the $m\le12$ sparse layers. The residual proof obligation is booked,
and no closure language is used for this layer.

\paragraph{The recorded dead zone.} The $42$ recorded dead-zone instances
of the base domain are discharged at \emph{witness} level: each admits a
machine-verified sharp decomposition with $\beta\le\lceil n/2\rceil$
(equality in $41$ of $42$; the sole strict case has $\beta=3<4$),
attributed as $13$ two-probe artifacts, $23$ cut-vertex witnesses and $6$
tight non-cut instances. This is a finite-instance disposition and does
not eliminate the corresponding premise of the conditional chain, which
remains open as stated (Section~\ref{int:subsec-conditional}).

\paragraph{Corrected paper-level values of the extended account.} The
extended-domain verification records corrected values, which supersede the
first-pass report: $1{,}279$ witness certificates ($484+795$; Section~\ref{com:subsec-kline}); the tight-layer validation
splits as $147$ even-$n$ instances for the tight-layer corollary proper
and $2{,}410$ odd-$n$ instances for its companion characterization
(the literal odd-$n$ proposition has $216$ counterexamples and is
superseded by the companion form ``success $\iff M(u)=d$ or
$M(u-1)\ge d-2$''); the insertion-depth account volume is $175{,}244$
under the corrected depth formula; and in the base domain's $d=6$
slice, $157$ instances are rescued ($156$ at insertion depth $2$, $1$
at depth $3$), zero otherwise.

\paragraph{The $K_8@28$ registered prediction.} For the placeability
programme: the phrase ``surviving regimes only $\{q=0,q=7\}$'' at
$K_8@28$ is a \emph{registered prediction}, not a measurement
(Remark~\ref{rem:claregistered}); the actual census found zero surviving
regimes on every line.

\paragraph{The Lov\'asz 1968 external step.} The all-odd support consumed
by Theorem~\ref{int:thm-Kgallai} rests on Theorem~\ref{int:thm-lovasz},
verified at statement level by five concordant published sources and
carried by the entire subsequent literature. Its offline proof
reconstruction was, however, judged \emph{not achieved}: the skeleton's
deletion step inflates the path count and its re-merging is blocked by
shared internal vertices, a failure pinned mechanically by $22$ witness
graphs on $n\le5$. Two transcription blemishes are on record alongside
(the equality reading ``$|D|=\lfloor n/2\rfloor$'' must be ``at
most''; the ``$\le n-1$ paths'' corollary is loose for even $n$), and
the consumed upper-bound-inflation line is cited in its corrected
upper-bound form, not as an identity. Page-by-page reading of the 1968
original remains mandatory before the external flag is lifted
(Section~\ref{int:subsec-conditional}).

\subsubsection{The generic-account limitations}\label{sta:subsubsec-genres}

The key-table ledger of Section~\ref{com:subsec-ledger} carries its four
conventions as limitations: the $N=14$ affine-normalized domain
equivalence is unproven; key-weighted and instance-weighted accountings
are separate; the no-$V_7$ reading is a one-way superset; and the
machine face of the omitted branch is frozen (its legacy
candidate-generation sources are lost, with the residual mass precisely
measured at $340$ positions for $N=8$ and $2{,}267{,}655$ at $N=14$).
Two further accounting rules bind quotations: per-instance accounts at
$N=8,9$ are listed separately from the key accounts, and the two
accounting layers of the falsification experiments ($12/900$ failing
instances versus $54/1{,}501$ failing witnesses;
Section~\ref{subsubsec:criteria}) are never interchanged. No ledger
result is extrapolated to $N\ge15$ or read across the two domain
conventions.

\subsubsection{The classification-programme boundaries}\label{sta:subsubsec-clares}

The placeability classification programme
(Section~\ref{subsec:echar}) is bounded as follows. The evidence
structure of the exception-set candidate \eqref{eq:claechar} is closed to
$k\le16$ on the empty-side and to $n=13$ on the tower side, with the
$k=17$ full-type sweep at $53{,}394{,}755$ types; the layers $k\ge18$
and the uniqueness side (closed only at $k=15$, refuted as a literal
statement at $k=21$) are open (Remark~\ref{rem:claregistry}). The
prime-family theorem covers the non-infinity regimes only; the infinity
segment of the prime family is open for every member
(Remark~\ref{rem:claprimscope}). The corrected Fourier identity
(Theorem~\ref{thm:clafourier}) is the only form of the character equation
cited anywhere in this paper, and all machine results of the programme
are independent of it (Remark~\ref{rem:clafourierfix}).

\subsubsection{The existence half of the criterion}\label{sta:subsubsec-gapex}

The general-$N$ existence half of the valuation criterion in its so1
form (Conjecture~\ref{int:conj-gapex}) is open: the per-$N$ closures at
$N\in\{6,8,10,12\}$ are theorem-form, the general statement is not, and
no sentence of this paper states otherwise. Three empirical obstacles are
registered against it --- the $\lambda=1$ dead corners ($2{,}528$
instances at $N=10$ with no one-leak witness), the failed-completion
triple $31/1{,}094/13{,}488$, and four falsified discriminators --- and
the candidate mechanisms (near-saturation normal forms; the leak-repair
channel; per-line board normal forms) each carry its own falsification
experiment. The associated $\mu$ $k=2$ stratified existence wall is open
on its paper face (machine face closed on seven domains; no logical
dependency on the family-image characterization;
Remark~\ref{rem:clabridge}).

\subsubsection{The GEN existence programme boundary}\label{sta:subsubsec-flicker}

The GEN existence programme carries one further measured boundary. Its
zero-dodge complete-walk counts on the $w_1=r-5$ flicker line are the
six exact values
\[
158,\;1{,}733,\;11{,}466,\;81{,}660,\;484{,}578,\;3{,}670{,}846
\]
for $r=19,21,23,25,27,29$ (the $r=29$ value exact but
class-bound-conditioned, quoted with that caveat) --- six tiers, not a
closed sequence. One further exact single point, $23{,}933{,}468$ at
$r=31$, stands on its own at exhaustive grade; at $r=33$ there is
\emph{no} count (a quantified hard wall; a segmented long run was started
and reports none). The even-$r$ blocking companion clears the
zero-dodge side only: complete walks with a positive zero-defect exist at
$r=20,22$ in the record and are untouched by it. The existence theorem
itself (registry grade, $r\ge43$) is supported by paper certificates and
step-feasibility, with its machine anchors corroborating only.

\subsection{Open problems}\label{sta:subsec-open}

The open problems below supersede the numbering of any earlier working
list; each is stated with what the present account gives and what it
withholds.

\begin{enumerate}[label=\textbf{O\arabic*.},leftmargin=2.4em]
\item \textbf{The remaining even columns of Case~4}
(Conjecture~\ref{conj:con27A}). Is
$\pn(K_{n_1,n_2})=\lceil n_1/2\rceil$ for all odd $n_1$ and even $n_2$,
$2\le n_2<n_1$? The ring-closure theorem
(Theorem~\ref{thm:conringclosure}) reduces each column to $h+1$ base
instances; the columns $n_2=2,4,6$ are closed, and five named bases are
open: $K_{15,8}$, $K_{17,8}$ (column $n_2=8$) and $K_{17,10}$,
$K_{19,10}$, $K_{21,10}$ (column $n_2=10$). The recorded negatives on
these bases are truncations, not exclusions
(Observation~\ref{obs:concostwall}): the decompositions may lie deeper in
the search tree, or the conjecture may fail. The next methodological step
suggested by the data is a constraint-propagation search driven by the
deficit ledger of Theorem~\ref{thm:constructgen} and the endpoint parity,
rather than deeper node budgets.
\item \textbf{Removing the structural conditions of the Fan-type
theorems} (Conjecture~\ref{conj:claAstar}). The alignment and separation
hypotheses fail exactly at \emph{remote high vertices}; by
Proposition~\ref{prop:claresidue}, a minimal counterexample must have a
bad partner at every high vertex. The remaining corner of the two-high
layer is blocked by the excess bound (zero violations on all $235$
computed instances), and the residue of the conjecture is that corner
together with graphs of at least three high vertices that are neither
separated nor covered by the profile parity. Note that the band-free
phenomenon is family-specific: outside the constructed family, the
general band-free statement is refuted by certified counterexamples
(Section~\ref{subsec:astar}).
\item \textbf{The even-$j$ Walecki residual.} Theorem~\ref{thm:conL0} is
for odd $j$; for even $j$ (blocks of order $1\bmod4$), the chain
construction would extend to all odd-order complete-block chains. In-chain
existence is machine-verified at the block values $K_5$ and $K_9$, but no
general construction is known. The natural route is a Hamilton
decomposition of $K_{2j+1}-C^*$ for a suitable cycle $C^*$ through the
prescribed path, which touches the Oberwolfach problem for odd orders.
\item \textbf{The general-$N$ existence half of the criterion}
(Conjecture~\ref{int:conj-gapex}). For all $N$ and all so1 instances with
capacity slack $\beta\ge2$, is there an admissible pair with
$\mathrm{val}\le1$? The $C_1$ side is closed for all $N$ by the
constructive one-leak theorem; the obstruction is $C_2$-side
reachability, with the three registered empirical obstacles and the
candidate mechanisms of Section~\ref{sta:subsubsec-gapex}, each carrying
its falsification experiment.
\item \textbf{The $\mu$ $k=2$ stratified existence wall.} The paper core
is the existence of $C_1$ with $q\le d_3+1$ and realized. The machine
face is closed on seven domains for the $k=2$ layer (zero violations);
the paper face has a recorded obstruction (no base-to-step mechanism);
and the problem has no logical dependency on the family-image
characterization (Remark~\ref{rem:clabridge}).
\item \textbf{The degree-$\ge4$ layer of the budget conjecture
$K$} (Definition~\ref{int:def-budget}). Beyond degree~$2$ the
deletion--reinsertion scheme meets its true wall; the computational
record shows the wall is an artefact of the method, not of $K$
(Section~\ref{sta:subsubsec-kres}). What is needed is a construction that
produces a $\lceil n/2\rceil$-decomposition with $D(x)\ge2$ directly,
without deleting $x$.
\item \textbf{The exception-set programme.} Beyond the closed faces of
Remark~\ref{rem:claregistry}: the layers $k\ge18$; uniqueness of the
minimal core for $k\ge21$; the prime-family infinity segment
(Remark~\ref{rem:claprimscope}); and the counting-side target of the
$\alpha''$ bridge (one-way and rigidity-free, machine base eight odd
points, $p\ge19$ unenumerated; Section~\ref{subsec:echar}).
\item \textbf{The $4$-degenerate residual.} The theorem of
\cite{ZLH24} covers $3$-degenerate graphs; the class of
$4$-degenerate graphs is, as far as we know, entirely open.
$K_{4,n}$ shows the interesting behaviour starts immediately: it is
$4$-degenerate and, by Theorem~\ref{thm:conk4odd}, beats its published
bounds by one.
\end{enumerate}

\subsection{Concluding remarks}\label{sta:subsec-conclusion}

\paragraph{What is proved.} On the constructive side: the exact value
$\pn(K_{4,n})=(n+1)/2$ for odd $n\ge3$ (two of whose instances are
formalized circumstances of one Lean theorem), the full $n_2=6$ column
and the structural reduction of the remaining even columns of Case~4 of
\cite{CE18}; the odd clique chain theorem
(Theorem~\ref{thm:conchain}) at its all-odd-$j$ domain; and the Fan-type
class theorems admitting a block of $\EV(G)$ of maximum degree at
least~$4$ (Theorems~\ref{thm:claAstar} and~\ref{thm:lasep}), at paper
level with formalization in progress. On the conditional side: the
budget framework, the implication $K\Rightarrow$ Gallai
(Theorem~\ref{int:thm-Kgallai}), the conditional chain, and the
descriptive registration of the eight-item premise ledger
(Theorem~\ref{int:thm-ledger}).

\paragraph{What is conditional.} The conditional route is
restated one last time in its honest form: the criterion's first
conjunct is a machine-closed candidate on explicitly tested domains with
a paper status open; every one of the eight premise conjuncts is an open
conjecture, an unverified external step, or a paper-open item with
machine backing; and no premise of the ledger is claimed, here or
elsewhere, to be within reach of the machine layer except as recorded.
The Lean layer sharpens checkability, not the claim surface: its
\texttt{sorryAx}-free theorems are those of
Table~\ref{sta:tab-leanitems}, its \texttt{native\_decide} certificates
are quoted with the kernel-reduction disclosure, and everything outside
the axiom probe is explicitly not claimed formalized
(Section~\ref{com:subsec-lean}).

\paragraph{What is measured.} The measured layer is reported at its
registered strengths throughout: the five exhaustively enumerated domains
of the budget conjecture (Table~\ref{com:tab-kdomains}; zero violations,
$208$ tool-dead instances all satisfying the conjecture); the nine-domain
key ledger under its four conventions
(Section~\ref{com:subsec-ledger}); the per-$N$ criterion closures with
zero residue on $588/30{,}760/634{,}960/7{,}798{,}903$ instances at
$N=6,8,10,12$; and the classification programme's four complete-graph
layers at their stated tiers (Section~\ref{subsubsec:firstlayers}). The
overall conjecture --- Gallai's path decomposition conjecture --- remains
open. What this paper provides is the verified portion of the landscape
around it, a conditional route whose every edge is checkable, and a
measurement apparatus whose every citation carries its reading
conventions (the four ledger conventions of Section~\ref{com:subsec-ledger}
and the tier qualifiers of Sections~\ref{sec:computation}
and~\ref{sta:subsec-leanitems}), so that the paper remains checkable
against exactly the evidence it quotes.

\paragraph{The no-claim constraints, restated.} The six no-claim
constraints of Section~\ref{int:subsec-claims} bind this closing section
as they bind the whole account; in the order stated there, and matching
the introduction's wording in content: no
closure at $L\ge13$ (candidate level only); no general-$N$ paper closure
of the covering identity (the existence half open); no completeness
claim for the four-piece (E)+(R) landing; no paper closure of the
general-$L$ uniform-construction account (the capacity-starved branch of
the channel-or-menu dichotomy); no full closure of that dichotomy
programme; and no
unconditional theorem-form statement of the valuation criterion --- its
Lean formalization is a bound under an unproved hypothesis, and a
criterion verified over finitely many $N$ is not an existence proof for
general $N$.

% end sec-status.tex
